\subsection{A Q-system realization of \texorpdfstring{$\sha{R8_3_721}$}{R8-3-721}}
\label{sec:r721-qsystem-model}
The earlier center calculations left this ring undecided. We now construct
an actual unitary fusion category and identify its Grothendieck ring within
the supplied complete census. The existence argument uses an inherited
Q-system; agreement of modular data is not used to manufacture an algebra.

\begin{theorem}[A cyclic generalized-Haagerup model]
\label{thm:r721-qsystem-model}
Let $\mathcal C$ be the $C_2$ de-equivariantization of the unitary
$C_8$ generalized-Haagerup category in
\cite[Example~4.4 and Section~6]{GrossmanIzumiCentersNext}.
Let $A=\one\oplus\rho$ be the image of its Q-system from
\cite[equations~(2.21)--(2.22)]{GrossmanIzumiCentersNext}.
Then $\mathcal D={}_A\mathcal C_A$ is a unitary fusion category of rank eight,
with trivial group of invertibles and FP dimensions
\[
 1,a,a,a+1,a+1,2a+1,2a,2a+2,
 \qquad a=\frac{3+\sqrt{17}}2.
\]
All its fusion coefficients are at most three. The complete supplied
rank-eight census therefore identifies
$K_0(\mathcal D)\cong\sha{R8_3_721}$.
\end{theorem}

We separate the categorical construction, the finite coefficient bound,
and the census identification. The last step uses the completeness of the
supplied census; it is not a new independent enumeration of fusion rings.

\begin{lemma}[An inherited-vertex obstruction]
\label{lem:cyclic-gh-invertibles-next}
Suppose an existing unitary generalized-Haagerup category for $C_{4m}$
satisfies the de-equivariantization hypotheses of
\cite[Section~6]{GrossmanIzumiCentersNext}, including
$\epsilon_{2m}(h)=(-1)^h$, and has the Q-system $\one\oplus\rho$
whose nonzero $\rho\to\rho^2$ vertex is proportional to $t_0$.
In its $C_2$ de-equivariantization, the Morita dual of this inherited algebra
has no nontrivial invertible object.
\end{lemma}
\begin{proof}
Write $g$ for the generator of the quotient group $C_{2m}$ and
$d=2m+\sqrt{4m^2+1}$. An invertible bimodule has underlying FP dimension
$1+d$. Since the source simple dimensions are $1$ and the irrational number
$d$, its underlying object has one invertible and one noninvertible
summand. If its invertible summand is $g^j$, free-module adjunction and
 equality of dimensions identify its underlying left and right modules with
$A g^j$ and $g^j A$, respectively. Their noninvertible summands are
$g^{-j}\rho$ and $g^j\rho$, so $2j=0$ in $C_{2m}$.

For $j=0$, the bimodule is an automorphism-twisted copy of $A$.
A unit-preserving algebra automorphism acts by a scalar $z$ on $\rho$.
The nonzero $\rho^2\to\one$ and $\rho^2\to\rho$ multiplication components
impose $z^2=1$ and $z^2=z$, hence $z=1$.

For $j=m$, use the crossed-product realization of
\cite[Section~6]{GrossmanIzumiCentersNext}. Its implementing unitary
$\lambda$ satisfies
\[
 \lambda^2=1,\qquad \rho'(\lambda)=\lambda,\qquad
 \alpha'_{2m}=\operatorname{Ad}(\lambda),\qquad
 \alpha'_m(t_0)=t_{2m}.
\]
Conjugation by $\alpha'_m$ sends $\rho'$ to
$\operatorname{Ad}(\lambda)\rho'$. The intertwiner
$u=\lambda:\rho'\to\operatorname{Ad}(\lambda)\rho'$ has
$u\otimes u=1$. Transporting the conjugate algebra to the object
$\one\oplus\rho'$ changes its vertex line from $\C t_0$ to
$\C t_{2m}\lambda$. These lines are orthogonal by the Cuntz relations.
A unit-preserving object isomorphism acts scalarly on the simple
$\rho'$ summand and cannot identify the two vertex lines. Thus
$g^m A g^{-m}\not\cong A$ as algebras. An invertible bimodule with underlying
free module $g^m A$ would give precisely such an algebra isomorphism,
contradicting this calculation.
\end{proof}
This lemma is conditional on existence of the stated generalized-Haagerup
category and Q-system. It makes no existence assertion for arbitrary $m$.

\begin{proof}[Proof of Theorem~\ref{thm:r721-qsystem-model}]
In the case $m=2$, the source category and Q-system are supplied by the
cited construction. The strong tensor de-equivariantization functor
preserves the Q-system, and its image $\rho$ is simple. Hence
$A=\one\oplus\rho$ is connected and separable, and its bimodule category
exists and is unitary. The source fusion rules are
\[
 g^4=1,\qquad g\rho=\rho g^{-1},\qquad
 \rho^2=\one+2\sum_{i=0}^3g^i\rho,
 \qquad \FPdim(\rho)=4+\sqrt{17}.
\]
Morita equivalence gives
$\FPdim(\mathcal D)=136+32\sqrt{17}$.

For free bimodule induction $I(X)=A\otimes X\otimes A$, adjunction gives
\[
 K_{ij}=\dim\operatorname{Hom}_{\mathcal D}(I(X_i),I(X_j))
       =[A X_j A:X_i].
\]
Every simple bimodule occurs in some $I(X)$, since its action map from the
free bimodule on its underlying object is a split epimorphism.
The known tensor-unit row is $(1,0,0,0,1,0,0,0)$ in source order
$(1,g,g^2,g^3,\rho,g\rho,g^2\rho,g^3\rho)$.
The remaining Gram matrix has rank four. Its complete finite factorization
checks $40$ pivot vectors and $14256$ multiplicity choices; it yields seven
row types and exactly four factors, of ranks $6,8,9,11$.

The known center in
\cite[Theorem~6.8 and Example~6.10]{GrossmanIzumiCentersNext}, equivalently
\cite[Definition~6.1 and Section~6.1.6]{GrossmanIzumiMetricNext}, removes
ranks nine and eleven. Indeed its twist-one support consists of six objects,
and the canonical Lagrangian multiplicity vector has form
$(1,p,x,y,z,w)$ on that support. The equations $S\ell=\ell$ and
$T\ell=\ell$ force
\[
 p=1,\qquad x+y=2,\qquad z+w=2.
\]
All nine nonnegative integral choices satisfy the full $S$ equations, and
\[
 \rank(\mathcal D)=\langle\ell,\ell\rangle
       =2+x^2+y^2+z^2+w^2\in\{6,8,10\}.
\]
Thus only the rank-six and rank-eight Gram factors remain.
The rank-six factor has an extra invertible bimodule with underlying object
$g^2+g^2\rho=g^2A$. Lemma~\ref{lem:cyclic-gh-invertibles-next} excludes it.
The actual dual therefore has the unique rank-eight factor, and dividing
its underlying dimensions by $\FPdim(A)$ gives the dimensions in the theorem.

It remains to bound its unknown fusion coefficients, without assuming a
bound in the construction. Let $B_{i x}=[U(x):X_i]$ be the unique rank-eight
forgetful matrix. Free induction satisfies
\[
 I(X)\otimes_A I(Y)\cong I(X\otimes A\otimes Y).
\]
For every source pair $(i,j)$ and dual output $z$, this gives
\[
 \sum_{x,y}B_{i x}B_{j y}N_{xy}^{z}
   =\sum_{a\in\{1,\rho\}}\sum_{u,t}
        N^{\mathcal C}_{i a}{}^{u}
        N^{\mathcal C}_{u j}{}^{t} B_{t z}.
\]
Together with the unit, duality, Frobenius reciprocity and both rational
components of the dimension character, these are $832$ exact linear
equations before redundant rows are removed. There are four possible
dualities on the two equal-dimension pairs; they give respectively
$120,112,112,104$ coefficient orbits. All four are retained.

Write one reduced system as $E n=b$, where $n\ge0$ consists of those fusion
coefficients. For every coordinate $j$, the supplied rational covector
$y_j$ satisfies
\[
 y_j E\ge e_j^{\mathsf T},\qquad y_j b\le3.
\]
Therefore $n_j\le y_j E n=y_j b\le3$. The certificates check every inequality
in exact rational arithmetic for all four dualities; numerical optimization
is used only to propose the covectors. This proves the required uniform
multiplicity bound for the actual category.

Finally, every requested dimension vector has sum $21+5\sqrt{17}$.
For a census tensor put $M_{jk}=\sum_iN_{ij}^{k}$. Its dimension vector is an
eigenvector of $M$ with that eigenvalue, so
$x^2-42x+16$ divides the integer characteristic polynomial of $M$.
An exact check of all $4500$ cumulative records leaves only
$\sha{R8_3_2906}$ and $\sha{R8_3_721}$.
Exact normalized eigenvectors in $\Q(\sqrt{17})$ give the wrong dimension
multiset for the former and the required multiset for the latter.
The unique match has source SHA-256
\par\noindent
\sha{76e26a815651476dfda5c2aa654f0d1327bdc890e08086cb6dbfe39ef9a1cdc5}.
Thus the constructed category has the based-isomorphism type
$\sha{R8_3_721}$ in the supplied complete census.
\end{proof}

\paragraph{Exact evidence and scope.}
The rank enumeration is
\path{evidence/rank8_next_20261021/semion_fermion/construction/r721_bimodule_rank.py}.
The coefficient system and rational certificates are in the same evidence
tree's \path{models/} directory:
\begin{quote}\small
\path{r721_dual_fusion.py}\quad and\quad
\path{r721_dual_multiplicity_bound.py}.
\end{quote}
The exhaustive crosswalk is
\path{evidence/rank8_next_20261021/r721/census_dimension_crosswalk.py}.
Normal and optimized-Python receipts, independent reconstructions, the
ordered tensor and the cited Q-system argument accompany these files.
The identification proves existence, and does not classify all associators
on this ring or assert that the Morita class is new.
